\documentclass[11pt]{article}

% ---------------------------------------------------------------------------
%  L6 — Perimeter gradings, king and square.
%  Category L (entirely machine-written); see paper/README.md and
%  docs/publication-split.md.
%
%  COMPUTE GATE CLEARED 2026-08-18.  Both k=6 censuses (n=78) finished on dalby
%  2026-08-09 and 08-10, 456/456 shards ok, binary git=6473890c, and were
%  harvested into results/ on 08-18.  §\ref{sec:k6} now carries the verdict:
%  onset 24, Phi_3^2 and no Phi_4 all CONFIRMED; the Phi_2 leading diagonal's
%  closed form REFUTED (5/2, not the predicted 15/4), and the k=5 identity of
%  the WHOLE Phi_3 block across lattices does not extend either.
%  The min-end runs were already done: square4 deep boxes (W=15,17) 08-08 and
%  ayr's square8 p=48 census 08-07.
%
%  NOVELTY STATUS: swept 2026-08-18.  The square-lattice column REPRODUCES
%  Asinowski-Barequet-Zheng, whose full text (ANALCO 2018) is now held and
%  read.  Do not claim the square column; the king column is a separate
%  derivation, NOT a corollary of their framework -- see sec:novelty.
%
%  Source material, all in this repository:
%    results/perimeter-defect-diagonals.md   the maximum end
%    results/perimeter-both-ends.md          both ends, the onset laws
%    docs/perimeter-both-ends-state.md       campaign state and the k=6 gate
%    results/min-site-perimeter.md           pmin on each lattice
%    cpp/perimeter_defect.cpp                the k = 2c + t identity, proved
% ---------------------------------------------------------------------------

% ---------------------------------------------------------------------------
%  paper/shared/preamble.tex — packages and macros common to every manuscript
%  in this directory.  Factored out of paper/polyplets-report.tex and
%  paper/technical-report.tex, which between them had two copies of the same
%  twenty lines.
%
%  technical-report.tex does NOT \input this file and is not to be edited to do
%  so: it is jasonp's prose and is read-only to the machine (paper/README.md,
%  "Who may edit what").  The duplication there is deliberate.
%
%  Assumes \documentclass[11pt]{article} has already been issued.
% ---------------------------------------------------------------------------
\usepackage[margin=1in]{geometry}
\usepackage{amsmath,amssymb}
\usepackage{amsthm}
\usepackage{booktabs}
\usepackage{graphicx}
\usepackage{numprint}
\usepackage{microtype}
\usepackage[table]{xcolor}
\usepackage[hidelinks]{hyperref}
\usepackage{flafter}   % a float never appears before the text that introduces it
\usepackage{float}     % [H] pins tables/figures exactly where written
\usepackage[font=small,labelfont=bf,labelsep=period]{caption}
\setlength{\intextsep}{16pt plus 3pt minus 3pt}

\npthousandsep{,}
\npfourdigitsep
\newcommand{\num}[1]{\numprint{#1}}
\newcommand{\g}{\cellcolor{gray!15}}

% Theorem environments.  One counter shared across kinds, so that a reference
% like "Theorem 4" is unambiguous in a paper that also has lemmas and
% propositions.
\theoremstyle{plain}
\newtheorem{theorem}{Theorem}
\newtheorem{proposition}[theorem]{Proposition}
\newtheorem{lemma}[theorem]{Lemma}
\newtheorem{corollary}[theorem]{Corollary}
\theoremstyle{definition}
\newtheorem{definition}[theorem]{Definition}
\newtheorem{example}[theorem]{Example}
\theoremstyle{remark}
\newtheorem{remark}[theorem]{Remark}
\newtheorem{conjecture}[theorem]{Conjecture}
\newtheorem{openproblem}[theorem]{Open Problem}

% Repo-path citations.  Every one of these must resolve in the tree that ships
% with the paper; `make gate-citations` checks the markdown, and
% paper/verify_claims.py checks the manuscripts.
\newcommand{\repo}[1]{\texttt{#1}}

% OEIS links.
\newcommand{\oeis}[1]{\href{https://oeis.org/#1}{#1}}

% Growth constants, so a rename is one edit.
\newcommand{\lam}{\lambda}
\newcommand{\muH}{\mu_H}

% ---------------------------------------------------------------------------
%  paper/shared/disclosure.tex — the two authorship disclosure blocks.
%
%  docs/publication-split.md §1 fixes the rule these implement:
%
%    "Under no circumstances will there be any ambiguity about whether a paper
%     is wholly my work (none), merely my prose (some) or entirely LLM-authored
%     (some)."   — jasonp, 2026-08-07
%
%  There are exactly two blocks and they live in exactly one file, so that a
%  paper cannot quietly soften its own disclosure.  The fixed sentences take no
%  arguments.  The ONE thing a paper supplies is the per-result verification
%  ledger, because that is the part that differs honestly between papers — and
%  §1 requires it to be per result, not in aggregate.
%
%  Do not add a \Pdisclosure or \Ldisclosure variant with an optional softening
%  argument.  If a paper does not fit one of these two blocks, the paper is
%  wrong, not the block.
% ---------------------------------------------------------------------------

\newsavebox{\disclosurebox}

\newenvironment{disclosureframe}
  {\begin{center}\begin{lrbox}{\disclosurebox}\begin{minipage}{0.93\textwidth}\small}
  {\end{minipage}\end{lrbox}\fbox{\usebox{\disclosurebox}}\end{center}}

% --- P: jasonp's prose -------------------------------------------------------
% Byline: Jason Parker, alone.  Argument: the per-result ledger of what the
% machine did.
\newcommand{\Pdisclosure}[1]{%
\begin{disclosureframe}
\textbf{Authorship.}\enspace Every sentence of this paper was written by its
named author. He has read every proof it states and believes each one, and he
is responsible for the correctness of what it claims.

\smallskip
\textbf{Machine contribution.}\enspace #1
\end{disclosureframe}}

% --- L: entirely machine-written --------------------------------------------
% Argument: the per-result verification ledger — for each result, what warrants
% it (Lean kernel, exact certificate, or labelled measurement) and how much of
% it a human has checked.
\newcommand{\Ldisclosure}[1]{%
\begin{disclosureframe}
\textbf{Authorship.}\enspace This document was written end to end by a large
language model. That includes its mathematics, its prose, its tables and its
\LaTeX{}. No sentence of it was written by a human. The mathematics it reports
was likewise derived by machine.

\smallskip
\textbf{Human role.}\enspace The person who directed this work chose what to
compute, chose what to publish, and is responsible for the decision to release
this document. Except where the ledger below says otherwise, that person has
not personally checked the arguments here, and this disclosure is not to be
read as an endorsement of them.

\smallskip
\textbf{What warrants each result.}\enspace #1
\end{disclosureframe}}

% --- Draft banner ------------------------------------------------------------
% For an L-paper whose verification ledger is not yet settled.  Loud on purpose:
% an L-paper without a truthful ledger must not be mistaken for a finished one.
%
% \firstdraftbanner is the standard wording, which five papers previously
% carried character-for-character in their own preambles.  It lives here for the
% same reason the disclosure blocks do: identical text in seven places drifts,
% and a paper must not be able to soften it by editing its own copy.  The
% optional argument is for a gate specific to one paper, appended verbatim; a
% paper needing different wording (L6, compute-gated) calls \draftbanner direct.
\newcommand{\firstdraftbanner}[1][]{%
\draftbanner{This is a first draft. The verification ledger below records what
warrants each result \emph{mechanically}; the column that records what a human
has read is empty, deliberately, because nobody has read it yet. Do not
circulate until that column says something true.#1}}

\newcommand{\draftbanner}[1]{%
\begin{center}
\fcolorbox{red!60!black}{yellow!15}{%
  \begin{minipage}{0.93\textwidth}\small
  \textbf{DRAFT — NOT FOR CIRCULATION.}\enspace #1
  \end{minipage}}
\end{center}}


\newcommand{\Z}{\mathbb{Z}}
\newcommand{\Q}{\mathbb{Q}}
\newcommand{\Phid}{\Phi}
\newcommand{\pmaxx}{p_{\max}}
\newcommand{\pminn}{p_{\min}}

\title{\textbf{Perimeter gradings of lattice animals}\\[4pt]
\large Two ends of one table, and two different onsets}
\author{[byline: jasonp's call --- \repo{docs/publication-split.md} \S1]}
\date{August 2026}

\begin{document}
\maketitle

\firstdraftbanner[ One further gate is specific to this paper. A literature
sweep ran on 2026-08-18 and the square-lattice column here \emph{reproduces}
published work now held in full (\S\ref{sec:novelty}); the king column and the
minimum end survived that sweep, but a negative from one pass is weaker than a
positive.]

\Ldisclosure{%
\emph{The identity $k = 2c + t$ and the monotonicity that licenses the search
prune} --- \textbf{proved}, and they are what make every measured value in this
paper trustworthy. The identity is \textbf{not ours}: it is
\cite{abz2018}'s $k = e + 2f$ for the square lattice, re-derived here for king
in ignorance of that (Remark~\ref{rem:abzidentity}).
\emph{The formulae of Tables~\ref{tab:maxend} and~\ref{tab:mincoeffs}} ---
\textbf{measured and interpolated, not proved}. Degrees and onsets are read off
data, not derived. Each row is an exact interpolation verified against at least
two spare points per residue class with the onset checked sharp, and each
predicted cyclotomic denominator is confirmed by exact division of the census
series --- at $k=5$ that is 28 consecutive exact zero coefficients the
interpolation never saw. That is strong evidence and it is not a proof.
\emph{The king min-end constants} --- measured against a predicted convolution
of the 4-coloured partition function at $i = 0..6$: fourteen of fourteen exact,
with both contributing mechanisms checked directly rather than inferred.
\emph{The tip series identification} --- $\phi_2$ of \cite{andrews1984}, an eta
quotient with an explicit infinite product, verified to $n = 40$. It is
\emph{not} a closed form.
\emph{$k = 6$} --- computed, and it decided four predictions three to one
(\S\ref{sec:k6}). The onset, the $\Phid_3$ exponent and the absence of
$\Phid_4$ hold; the $\Phid_2$ leading diagonal's closed form does not, and is
withdrawn here rather than repaired.
\emph{Novelty} --- pass run 2026-08-18, and it found that the defect
identity, the cyclotomic-denominator theorem and the degree column are all
\cite{abz2018}'s, the last as a conjecture this paper confirms.
\S\ref{sec:novelty} lists what survives and what does not.
\emph{Human verification:} none, as of this draft.}

\begin{abstract}
\noindent
Let $A(n,p)$ count fixed lattice animals of $n$ cells with site-perimeter $p$.
This table has two boundary curves, and grading from each gives a different kind
of object.

At the \textbf{maximum} end the extremal animals are sticks, the natural index
is the defect $k = \pmaxx(n) - p$, and the classes are quasi-polynomials in $n$
of degree exactly $k$. Their onsets follow
$\mathrm{onset}(k) = \binom{k+1}{2} + 3$, verified at five consecutive values on
each of the square and king lattices. Through $k = 6$ the period, the degree,
the onset \emph{and the leading coefficient} are identical on the two lattices;
only sub-leading coefficients differ. The denominators are cyclotomic, with
$\Phid_d$ first appearing at $k = 2d-1$, and the coefficients form their own
triangle whose leading diagonal is lattice-independent and whose every lower
diagonal is not.

At the \textbf{minimum} end the extremal animals are balls of the adjacency
metric, $\pminn(n)$ grows like $\sqrt n$ and is a step function, so the natural
index is $p$ rather than $n$. There the classes are eventually \emph{constant},
with onset linear in the defect. On the king lattice the constants are a
convolution of the 4-coloured partition function with the box-skew deficits ---
fourteen of fourteen predicted values exact. What decides whether a column
stabilises at all is not parity but \emph{attainability}: a column stabilises
exactly when its $p$ is realised as $\pminn(n)$ for some $n$.

The mechanism behind the min-end constants is a Young diagram at each corner of
the isoperimetric hull, giving a factor $P(x)^c$ with $c$ the corner count ---
confirmed on a hexagonal hull as well as a square one. The square lattice's
diamond is an apparent exception: its tips are not lattice-aligned, and each
counts the order ideals of its own tangent cone, of which $P(x)$ is the aligned
special case. The diamond's tip series is Andrews' $\phi_2$, so the two corner
types here are the first two members of one family.

\end{abstract}

\tableofcontents

\section{Introduction}

The site perimeter of a lattice animal is the number of empty cells adjacent to
it. Write $A(n,p)$ for the number of fixed animals of $n$ cells with site
perimeter $p$. For each $n$ the possible $p$ occupy an interval, and the two
boundary curves of that interval differ in kind:
\[
  \pmaxx(n) = \begin{cases} 2n+2 & \text{square lattice} \\ 4n+4 & \text{king lattice}\end{cases}
  \qquad\text{is linear,}
\]
while $\pminn(n)$ grows like $\sqrt n$ and is a step function. So the two ends
want different indices, and this paper grades from both.

\paragraph{Which lattices.} Throughout, ``square'' means the rook-connected
animals (polyominoes) and ``king'' means the king-connected ones (polyplets).
Where a third lattice appears it is the six-neighbour triangular one, and it is
used only to test whether a mechanism generalises.

\section{The maximum end}
\label{sec:max}

\subsection{The defect is monotone}

Grading by $k := \pmaxx(n) - p$ looks as though it should want a
surplus-budgeted row transfer, on the model of the height grading in companion
paper L1. It cannot: the perimeter defect does not bound the transverse extent
of a linear sweep, since a U whose arms are far apart has $k = 3$ and unbounded
spread. That is presumably why the published method classifies patterns instead
of sweeping.

Something simpler works. Let $c = e - n + 1$ be the cycle rank of the animal's
adjacency graph, and let $t$ be the sum, over empty cells adjacent to the
animal, of (number of animal neighbours $- 1$).

\begin{proposition}
\label{prop:kct}
$k = 2c + t$, with both terms nonnegative. Consequently $k$ never decreases when
a cell is added to an animal.
\end{proposition}

\begin{remark}[this identity is published, for the square lattice]
\label{rem:abzidentity}
Proposition~\ref{prop:kct} is not ours on the square lattice. Asinowski,
Barequet and Zheng~\cite{abz2018} state it as $k = e + 2f$, with $e$ the total
excess summed over perimeter cells and $f$ the number of inner faces of the dual
map --- the circuit rank. That is the present statement with $c = f$ and
$t = e$, and their proof is the same accounting: four neighbours per cell, less
the excesses, less two per spanning-tree edge, less two per cycle-closing edge.
It was re-derived here for the king lattice without knowing that, and the
king case does not follow from theirs by substitution, since the per-cell degree
and hence $\pmaxx$ changes, but the identity and its proof idea are theirs.
Their full text was obtained after this paper was drafted, and it confirms the
reading: their accounting is $p = 6n - e - 2|E|$ in three dimensions, the same
two-sources-of-loss argument, with $|E| \ge n-1$ giving the bound and the
circuit rank giving the excess over a spanning tree.
\end{remark}

\begin{proof}[Sketch, with the full argument in the enumerator's header]
Placing a cell raises $\pmaxx$ by $\deg/2$ and changes $p$ by $-1 + g$, where $g$
counts the neighbours that become fresh perimeter cells; $g$ is capped by the
neighbours the new cell shares with the animal cell it touches. The bound is
tight on both lattices --- a stick extends at $\Delta k = 0$.
\end{proof}

So a Redelmeier depth-first search that abandons a partial animal the moment
$k > k_{\max}$ is \emph{exact}. A second monotonicity makes it fast: a candidate
cell's cost $\Delta k$ is itself nondecreasing as the animal grows, so a
candidate over budget is over budget forever and can be dropped from the untried
list rather than re-tested at every node. That is worth about $20\times$ --- at
$n = 20$, $k \le 5$ on the king lattice, $2.4\times10^8$ nodes fall to
$1.4\times10^7$.

The enumerator reaches $n = 70$ on both lattices on a single core ($913$~s on
square, $2017$~s on king), where brute force stops at $n = 14$ on king. It is
gated: pruned counts match an independent brute-force census cell for cell on
both lattices, the prune is confirmed to change nothing against an unpruned
control run, and the $(c,H)$ marginals sum back.

\subsection{The classes and their onsets}

\begin{table}[H]
\centering\small
\begin{tabular}{r r r r l l}
\toprule
$k$ & period & degree & onset & denominator & leading coefficient \\
& & & & & \textbf{(both lattices)} \\
\midrule
0 & 1 & 0 & 2  & $\Phid_1$ & 2 \\
1 & 1 & 1 & 3  & $\Phid_1^2$ & 4 \\
2 & 1 & 2 & 6  & $\Phid_1^3$ & 6 \\
3 & 2 & 3 & 9  & $\Phid_1^4\Phid_2^2$ & $13/2$ \\
4 & 2 & 4 & 13 & $\Phid_1^5\Phid_2^3$ & $17/3$ \\
5 & 6 & 5 & 18 & $\Phid_1^6\Phid_2^4\Phid_3$ & $593/144$ \\
6 & 6 & 6 & 24 & $\Phid_1^7\Phid_2^5\Phid_3^2$ & $13325/5184$ \\
\bottomrule
\end{tabular}
\caption{The defect-$k$ classes. Period, degree, onset and leading coefficient
agree on the square and king lattices at every $k \le 6$; sub-leading
coefficients do not.}
\label{tab:maxend}
\end{table}

\begin{quote}
\textbf{Onset.} $\mathrm{onset}(k) = \tbinom{k+1}{2} + 3$ for $k \ge 2$.
\end{quote}

\begin{quote}
\textbf{Degree.} $\deg Q_k = k$, equivalently $\Phid_1^{k+1}$ exactly, for
every $k \le 6$ on both lattices.
\end{quote}

That row is not an observation of ours either. \cite{abz2018} conjecture, in
any dimension and for any fixed defect, that the highest-degree factor of the
characteristic polynomial is $(x-1)^{k+1}$, contributed by the pattern of
$k+1$ independent legs, so that the count is asymptotically $\gamma n^k$.
The square column here is evidence for that conjecture at $d = 2$ through
$k = 6$, and the king column is evidence for the same statement on a lattice
their framework does not reach.

Five consecutive hits on each lattice ($6, 9, 13, 18, 24$). The degenerate
cases $k = 0, 1$ have onsets $2$ and $3$ against the formula's $3$ and $4$,
which is exactly why scanning the whole sequence $2,3,6,9,13,18$ found nothing.
The $k = 6$ hit is the first that was predicted before it was measured, and it
was measured without interpolating: the smallest onset admitting \emph{any}
cyclotomic denominator in the division test of \S\ref{sec:k6} is exactly $24$
on both lattices. The formula predicts $\mathrm{onset}(7) = 31$.

$k = 3$ carries a genuine period-$2$ term on \emph{both}
lattices, so quasi-polynomiality is intrinsic to the perimeter grading and not a
square-lattice artefact. On the square lattice, for $n \ge 9$,
\[
  Q_3(n) = \tfrac{13}{2}n^3 - 89n^2 + \tfrac{1947}{4}n - \tfrac{2107}{2}
           + (-1)^n\Bigl(\tfrac{5n}{4} - \tfrac{21}{2}\Bigr),
\]
the parity part having degree $1$, not $0$. On the king lattice, same shape,
different sub-leading coefficients:
\[
  Q_3(n) = \begin{cases}
    \tfrac{13}{2}n^3 - 69n^2 + 360n - 860 & n \text{ even},\\[2pt]
    \tfrac{13}{2}n^3 - 69n^2 + \tfrac{715}{2}n - 839 & n \text{ odd}.
  \end{cases}
\]

\subsection{The cyclotomic content}
\label{sec:cyclo}

At $k = 5$ the fitted period is $6$, but the $\Phid_6$ component of the constant
term is \emph{exactly zero}: what looks like period $6$ is $\Phid_2$ and
$\Phid_3$ acting independently, not a primitive sixth root. Reading exponents off
where each factor first perturbs a coefficient ($\Phid_2$ enters the $n^3$
coefficient, $\Phid_3$ only the constant), gives, across $k \le 6$,
\[
  D_k \;=\; \Phid_1^{\,k+1}\,\Phid_2^{\,k-1}\,\Phid_3^{\,k-4},
  \qquad\text{with}\qquad
  \boxed{\;\Phid_d \text{ first appears at } k = 2d-1\;}
\]
($d=2$ at $k=3$, $d=3$ at $k=5$; $\Phid_4$ is correspondingly absent at
$k = 6$, which was tested).

That the denominators are cyclotomic at all is, again, published on the square
lattice: \cite{abz2018} prove that for each fixed $k$ the generating function
of $(A(n, 2n+2-k))$ is rational \emph{with denominator a product of cyclotomic
polynomials}, and record the same main result in $d$ dimensions for polycubes.
So the frame of this subsection is the king analogue of their theorem. It is an
analogue and not a corollary: their proof shrinks maximal runs of grid slices
whose projection is a set of pairwise non-adjacent cells, to a unique reduced
representative per pattern class, and bounds the number of such classes through
excess cells, $L$-cells and degree-1 cells. Every one of those steps is stated
under face connectivity --- the perimeter accounting, the definition of an
$L$-cell as an occupied cell with two occupied neighbours that are not
opposite, the non-adjacency condition defining a cut, and a handshake bound at
maximum degree $6$. King adjacency breaks the cut condition in particular. The
shape of the argument transports; the theorem does not, which is why
the king column of \S\ref{sec:max} is measured rather than cited. What is
not in their work, as far as the pass of
\repo{docs/priority-passes-2026-08-18.md} reached, is the exponent rule boxed
below, which $\Phid_d$ appears at which $k$, and everything the
coefficient triangle carries. If that rule holds, ``quasi-polynomial'' is a
coarse label and each successive prime-power periodicity costs a fixed two
units of defect.

\paragraph{The denominators are predicted, then checked.} They are not fitted.
$D_k$ is read off the table above and the census series (through $n = 70$ for
$k \le 5$, $n = 78$ for $k = 6$) is divided by it; a polynomial must be left,
and one is, for every $k$ on both lattices, with $\gcd(N,D) = 1$ so no factor is
spurious. At $k = 5$ that is \textbf{28 consecutive exact zero coefficients} the
interpolation never saw, and at $k = 6$ \textbf{39}: a far harder test than the
two-spare-points-per-class bar at which the formulae were fitted. Each passing
denominator is then retested with one factor removed at a time, so the
denominator reported is the minimal one.

\subsection{The coefficient triangle}
\label{sec:coefftri}

Writing the partial fractions with $R_k$ the polynomial part carrying the
pre-onset holdouts, the king generating functions are
\[
\begin{aligned}
  G_0 &= R_0 + \tfrac{2}{\Phid_1}, \\
  G_1 &= R_1 + \tfrac{4}{\Phid_1^2} - \tfrac{12}{\Phid_1}, \\
  G_2 &= R_2 + \tfrac{12}{\Phid_1^3} - \tfrac{48}{\Phid_1^2} + \tfrac{92}{\Phid_1}, \\
  G_3 &= R_3 + \tfrac{39}{\Phid_1^4} - \tfrac{216}{\Phid_1^3}
        + \tfrac{2445/4}{\Phid_1^2} - \tfrac{5135/4}{\Phid_1}
        + \tfrac{5/4}{\Phid_2^2} - \tfrac{47/4}{\Phid_2},
\end{aligned}
\]
and so on through $G_5$, whose $\Phid_3$ block is $(40x/27 + 8/3)/\Phid_3$. Each
block is one piece of the quasi-polynomial: $a_j/\Phid_1^j$ contributes
$a_j\binom{n+j-1}{j-1}$, the plain part; $b_j/\Phid_2^j$ the same times
$(-1)^n$; the $\Phid_3$ block a bounded period-3 wobble.

Those coefficients form their own triangle, and the lattice-independence in it
is \emph{exactly one diagonal deep}. Writing rows by $k$ and columns by the
offset $d$ from the top of each block, with $\ast$ marking entries identical on
square and king:

\begin{table}[H]
\centering\small
\begin{tabular}{r r r r r}
\toprule
$\Phid_1$ block & $d=0$ & $d=1$ & $d=2$ & $d=3$ \\
\midrule
$k=0$ & $2^\ast$ & & & \\
$k=1$ & $4^\ast$ & $-12^\ast$ & & \\
$k=2$ & $12^\ast$ & $-48$ & $92$ & \\
$k=3$ & $39^\ast$ & $-216$ & $2445/4$ & $-5135/4$ \\
$k=4$ & $136^\ast$ & $-1911/2$ & $14049/4$ & $-18863/2$ \\
$k=5$ & $2965/6^\ast$ & $-4234$ & $685645/36$ & $-2225243/36$ \\
$k=6$ & $66625/36^\ast$ & $-168514/9$ & $21478331/216$ & $-27170587/72$ \\
\midrule
$\Phid_2$ block & $d=0$ & $d=1$ & $d=2$ & $d=3$ \\
\midrule
$k=3$ & $5/4^\ast$ & $-47/4^\ast$ & & \\
$k=4$ & $5/4^\ast$ & $-47/2$ & $349/2$ & \\
$k=5$ & $15/8^\ast$ & $-293/8^\ast$ & $10503/32$ & $-28141/16$ \\
$k=6$ & $5/2^\ast$ & $-1843/32$ & $37807/64$ & $-3556$ \\
\bottomrule
\end{tabular}
\caption{King values, truncated at $d = 3$. Every $d = 0$ entry agrees across
the lattices; $d = 1$ already fails from $k = 2$ in the $\Phid_1$ block, and the
$k = 5$ agreement at $d = 1$ in the $\Phid_2$ block does not survive to
$k = 6$ ($-1843/32$ against square's $-1831/32$).}
\label{tab:coefftri}
\end{table}

The statement is a diagonal one about a triangle of \emph{coefficients} rather
than about counts, which is the closest structural analogue to the
height-diagonal law that this grading offers.

Neither leading diagonal closes.

\begin{itemize}
\item \textbf{$\Phid_1$:} $a_{k+1} = c_k\,k! = 2, 4, 12, 39, 136,
  593\cdot120/144, 66625/36$ --- integral through $k = 4$ and then not. No
  closed form.
\item \textbf{$\Phid_2$:} $5/4$, $5/4$, $15/8$ fits $b_{k-1} =
  5(k-2)!/2^{\,k-1}$, which predicts $15/4$ at $k = 6$. The measured value is
  \textbf{$5/2$}, on both lattices. Three terms, one formula, and it is wrong at
  the fourth; the entry above is a coincidence and is withdrawn. What survives
  is the weaker and more robust statement, that the value is the same on both
  lattices, which is the universality claim and not a closed form.
\end{itemize}

Two identities that looked exact at $k = 5$ were the same kind of accident. The
$\Phid_3$ block at $k = 5$ is identical on both lattices \emph{in full},
$(40x/27 + 8/3)/\Phid_3$, which invited ``the entire period-$3$ content is
lattice-independent''. At $k = 6$ only its \emph{leading} coefficient
($-5920/243$) is shared: king carries
$(-5920x^3 - 19128x^2 - 18528x - 12656)/243$ against square's
$(-5920x^3 - 18768x^2 - 18168x - 12296)/243$. Likewise the $\Phid_2$ block's
$d = 1$ entry agrees at $k = 3$ and $k = 5$ and disagrees at $k = 6$.

The rule that survives all three corrections is the one this section opened
with, now tested one row further: \textbf{the leading diagonal of
each cyclotomic block is lattice-independent, and nothing below it is.}

\subsection{The recentred basis is nonnegative integers}

The partial-fraction basis $1/\Phid_1^j \mapsto \binom{n+j-1}{j-1}$ is centred at
$n = 0$, which is why its coefficients are the awkward rationals above.
Re-expand each residue class in $\binom{m}{j}$ with $m = (n - n_0)/\text{period}$
stepping along the class from its first in-regime point, and \emph{every
coefficient is a nonnegative integer}, on both lattices, in
every residue class, for $k = 2,\dots,5$ (the $k = 6$ recentring has not
been run):
\[
\begin{aligned}
  k = 2:&\quad 92,\,48,\,12 \quad(\text{king}); \qquad 60,\,40,\,12 \quad(\text{square}) \\
  k = 3,\ r = 1:&\quad 1528,\,1868,\,1164,\,312; \qquad 856,\,1324,\,1004,\,312 \\
  k = 4,\ r = 1:&\quad 45753,\,48873,\,32438,\,12484,\,2176; \qquad 22993,\,30293,\,24222,\,10956,\,2176
\end{aligned}
\]
The leading coefficient is lattice-independent ($312$, $2176$, $3842640$ for
$k = 3,4,5$) --- it is $c_k\,k!\,\text{period}^k$, so this is the known leading
diagonal in integer form. Defining the recentring at all needs
the onset formula, so this is downstream of it.

\subsection{Doors closed}
\label{sec:closed}

Each of these was an attractive reading, and each is dead.

\begin{itemize}
\item \textbf{The parity is not carried by the cyclic animals.} Splitting each
  class by cycle rank: the $c \ge 1$ parts are plain polynomials and tiny
  (square $k=3$, $c=1$ is the constant $8$; king $k=3$, $c=1$ is one animal at
  $n=4$ and nothing after). All of the $(-1)^n$ sits in the \emph{acyclic}
  animals, so the appealing ``a ring has an even cell count'' story is wrong.
\item \textbf{It is not carried by any single height slice.} Every slice
  $H = n - j$ at fixed $j$ is a plain polynomial on both lattices.
\item \textbf{It is not the height floor.} The floor identity of \S\ref{sec:identity}
  has period $k+1$, and the near-floor king slices do oscillate with period $4$
  at $k=3$ --- but it predicts period $5$ at $k=4$ where the answer is $2$, and
  the square lattice has no meaningful floor at all yet shows the identical
  period. The floor contributes; it is not the source.
\item \textbf{$c_k\,k! = 2\binom{2k}{k}$ is dead.} The products run
  $2, 4, 12, 39, 136$ against $2, 4, 12, 40, 140$ --- tempting through $k = 4$;
  at $k = 5$ the leading coefficient is $593/144$, whose product with $5!$ is
  not even an integer.
\item \textbf{The numerators are not nonnegative.} Done properly on the tail
  series, $\widetilde N_k$ is a polynomial of degree exactly $\deg D_k - 1$,
  itself a fresh confirmation of the predicted denominators, but its
  coefficients change sign from $k = 2$ on.
\item \textbf{The onset is not the positivity threshold.} The least $n_0$ at
  which all forward differences are nonnegative is strictly \emph{below} the
  onset in every class on both lattices ($k=5$: threshold $12$ against onset
  $18$). Both facts are real; neither explains the other.
\end{itemize}

\section{The minimum end}
\label{sec:min}

$\pminn(n)$ grows like $\sqrt n$ and is a step function, so an $n$-indexed ladder
there would have columns polynomial in $\sqrt n$ at best. The right row index is
$p$:
\[
  C(p,i) \;:=\; A\bigl(n_{\max}(p) - i,\; p\bigr),
  \qquad n_{\max}(p) = \max\{n : \pminn(n) \le p\}.
\]
$\pminn$ itself is not ours and is not refitted here: on the square lattice it is
\oeis{A261491}, $\lceil 2 + \sqrt{8n-4}\rceil$, and on the king lattice
\oeis{A235382}, $2\lceil 2\sqrt n\rceil + 4$. Both were re-verified against the
census before anything else ran.

Data at this end comes from \emph{complementation} rather than growth: enumerate
the isoperimetric hulls and remove cells, since the defect is not monotone under
addition here.

\subsection{King: the constants are the 4-coloured partition function}

$C(p,i)$ is eventually constant in $p$ along each residue class mod $4$, and the
constants are
\[
\begin{aligned}
  p \equiv 0 \pmod 4:&\quad C(i) = q_4(i) + 2\sum_{s\ge1} q_4\bigl(i - s^2\bigr), \\
  p \equiv 2 \pmod 4:&\quad C(i) = 2\sum_{s\ge0} q_4\bigl(i - s(s+1)\bigr), \\
  p \text{ odd}:&\quad 0,
\end{aligned}
\]
with $q_4(j) = [x^j]\prod_n (1-x^n)^{-4} = 1, 4, 14, 40, 105, 252, 574$ the
4-coloured partition numbers.

\begin{table}[H]
\centering\small
\begin{tabular}{l r r r r r r r}
\toprule
$i$ & 0 & 1 & 2 & 3 & 4 & 5 & 6 \\
\midrule
$p \equiv 0 \pmod 4$ & 1 & 6 & 22 & 68 & 187 & 470 & 1106 \\
$p \equiv 2 \pmod 4$ & 2 & 8 & 30 & 88 & 238 & 584 & 1360 \\
\bottomrule
\end{tabular}
\caption{Measured against predicted: all fourteen exact.}
\label{tab:mincoeffs}
\end{table}

The $p \equiv 0$ row is now confirmed at five separate perimeters. The census
reaches $p = 48$, and at $p = 32, 36, 40, 44, 48$ (bounding boxes of side $7$
through $11$) all seven columns read $1, 6, 22, 68, 187, 470, 1106$ with no
drift. Below that the columns are visibly still climbing, and each reaches its
value exactly at the predicted rung: $i = 0$ at $p = 8$, $i = 1$ at $12$,
$i = 2$ at $16$, and so on to $i = 6$ at $p = 32$. Seven of seven onsets land on
$4i + 8$.

Two mechanisms do the work, and both are checked directly rather than inferred.

\paragraph{Where $q_4$ comes from.} Removing a corner cell of a filled king box
drops one ring cell and adds the removed cell: net zero. So the
perimeter-preserving removals form a Young diagram at each of the four corners,
independently, giving $P(x)^4$. The per-box free-removal counts converge to
$1,4,14,40,105,252,574$ term by term, and converge \emph{exactly} when the box
side exceeds the removal count --- $W=5$ is right to $j=4$, $W=6$ to $j=5$,
$W=7$ to $j=6$.

\paragraph{Where the period 4 comes from.} At semiperimeter $S = w + h$ the
balanced box is optimal, and a skew of $s$ costs $s^2$ cells when $S$ is even and
$s(s+1)$ when odd. Different deficit sets, hence different columns, hence period
$4$ in $p$.

\begin{quote}
\textbf{Stabilisation onset.} $p^\ast = 4i + 8$ for $p \equiv 0$ and $4i + 10$
for $p \equiv 2 \pmod 4$, exact for $i = 0,\dots,6$.
\end{quote}

Both say the same thing, that the balanced box side must exceed $i$, and both
are \emph{linear} in the defect, against the maximum end's triangular
$\binom{k+1}{2}+3$.

\subsection{What decides stabilisation}

King's odd-$p$ rows do not stabilise; they grow. The square lattice's odd-$p$
rows do stabilise, so parity is not the rule. The rule, tested on both lattices,
is attainability:

\begin{quote}
$C(p,\cdot)$ stabilises $\iff$ $p$ is \textbf{attained} as $\pminn(n)$ for some
$n$.
\end{quote}

King's $\pminn = 2\lceil2\sqrt n\rceil + 4$ is always even, so no odd $p$ is ever
an isoperimetric perimeter, and every odd column counts animals that are one unit
worse than any optimum --- a different population, with a positional freedom
along the boundary that grows with it. The square lattice's
$\pminn = 2 + \lceil\sqrt{8n-4}\rceil$ attains every integer in range
\emph{except} $5$, and correspondingly every one of its four columns stabilises.
The hypothesis was checked class by class against the measured stabilisation on
both lattices, with no exceptions.

The non-attained columns are quadratic in $p$, fitted exactly with holdouts:
\[
  p \equiv 1\ (4),\ i{=}1: \tfrac{p^2-10p-39}{16};\qquad
  p \equiv 3\ (4),\ i{=}1: \tfrac{p^2-10p-43}{8};\qquad
  p \equiv 1\ (4),\ i{=}2: \tfrac{3p^2-30p-197}{8}.
\]

\subsection{The free-removal factor is \texorpdfstring{$P(x)^{\text{corners}}$}{P(x)\^{}corners}}
\label{sec:tips}

Measured directly on single large hulls, the perimeter-preserving removal counts
are
\[
\begin{aligned}
  \text{king, square box (4 corners), } W{=}7:&\quad 1, 4, 14, 40, 105, 252, 574 = [x^j]P(x)^4, \\
  \text{tri6, hexagon (6 corners), } r{=}4:&\quad 1, 6, 27, 98, 315 = [x^j]P(x)^6,
\end{aligned}
\]
with $P(x) = \prod_n (1-x^n)^{-1}$, each converging term by term as the hull
outgrows $j$. So the per-hull factor is $P(x)^c$ with $c$ the number of corners
of the isoperimetric hull --- a Young diagram at each corner, independently. The
hexagon confirms that this was a general argument and not a fact about boxes.

\paragraph{The square lattice's diamond is an apparent exception.} It gives
\[
  1, 4, 18, 60, 187, 524, 1388, 3452, 8229, 18800, 41536 ,
\]
not $P(x)^4$, and is larger from $j = 2$ on: $18$ against $14$. A fourth root gives a per-tip series
$1, 1, 3, 5, 9, 14, \dots$, not the partition numbers $1, 1, 2, 3, 5, 7$.

\paragraph{The tip is not lattice-aligned.} A diamond tip is $90^\circ$ off the square lattice's axes,
so the removable shapes there are not plain Young diagrams. The tip counts
the \textbf{order ideals of its tangent cone} $\{a \ge |b|\}$:
$1, 1, 3, 5, 9, 14, 24, 35, 55, 81, \dots$, whose fourth power reproduces every
measured diamond term. $P(x)$ is the lattice-aligned special case, so king, tri6
and square4 collapse to one statement and the diamond stops being an exception.

The last two diamond terms are a prediction that was tested after the fact. The
$W = 15$ and $W = 17$ hulls were enumerated once the identification was already
written down, and returned $3452$ at $j = 7$ and $8229$ at $j = 8$, which are
$[x^7]g^4$ and $[x^8]g^4$ exactly. The convergence rule shows its teeth here:
$W = 15$ gives $8193$ at $j = 8$, where $j = 8$ needs $W \ge 17$, and only the
wider hull reports the converged $8229$.

That last step rested on one hull, so it was checked against a second
implementation. Counting free removals by growing the set one cell at a time,
rather than by filtering every subset, is cheap enough to reach hulls the
enumerator cannot afford; it reproduces both enumerated rows exactly, and
returns $8229$ again at $W = 19$, $21$, $23$ and $25$. The same runs carry the
series to $18800$, $41536$ and $88924$, each at two hulls or more, and break at
$j = r+1$ at every radius --- the convergence rule holding at six more places
than it was checked at.

\paragraph{A refutation.} The six-term prefix
$1,1,3,5,9,14$ matches \oeis{A120452}, which continues $23, 34, 52, 75,\dots$
If the diamond's four tips were independent then $g^4$ with $g = $ A120452 would
give the square lattice's $j = 6$ free-removal term as $1384$. The measurement is
$1388$, and the convergence rule, that term $j$ is correct once $W = 2j+1$,
verified at $j = 4$ and $j = 5$, says the run was already converged; wider runs
at $W = 15$ and $W = 17$ confirm it. So A120452 is wrong at its seventh term:
$23$, where the correct value is $24$.

\paragraph{The tip series has a name.} It is \oeis{A053993}, Andrews'
$\phi_2$~\cite{andrews1984}: generalized Frobenius partitions with up to two
repetitions per row. So it is an \textbf{eta quotient} with an explicit infinite product (Andrews
eq.~5.9), verified here to $n = 40$: as explicit as $P(x)$ itself, and not a
closed form.

The reading is that the two corner types in this problem are the first two
members of one family: a unimodular corner gives $\phi_1 = P(x)$, and the
index-$2$ diamond tip gives $\phi_2$. The natural guess, $\phi_m$ for index
$m$, is \textbf{false} at $m = 3$, and is recorded here as a closed door rather
than left as a suggestion.

\subsection{The square lattice's min-end ladder}

All four residue classes stabilise here, the odd ones included.

\begin{table}[H]
\centering\small
\begin{tabular}{l r r r r r l}
\toprule
 & $i=0$ & $i=1$ & $i=2$ & $i=3$ & $i=4$ & stabilises from \\
\midrule
$p \equiv 0 \pmod 4$ & 1 & 9  & 52  & 206 & 719 & $p^\ast = 4i+4$ \\
$p \equiv 2 \pmod 4$ & 4 & 22 & 106 & 392 & --- & $p^\ast = 4i+6$ ($i\ge1$) \\
$p \equiv 1 \pmod 4$ & 4 & 28 & 124 & 456 & --- & $p^\ast = 4i+9$ \\
$p \equiv 3 \pmod 4$ & 4 & 28 & 124 & 456 & --- & $p^\ast = 4i+7$ \\
\bottomrule
\end{tabular}
\caption{The linear stabilisation onset holds here too, exactly, with a
different constant per class --- four or five points per class, no exceptions
above $i = 0$. The two odd classes agree with each other on every converged
value. $C(p,0) = 1$ when $p \equiv 0$ (the perfect diamond, unique) and $4$
otherwise (a partial layer in four rotationally equivalent positions), with
$p = 6$ the small-$n$ exception.}
\label{tab:sq4min}
\end{table}

\section{An exact identity relating \texorpdfstring{$n$, $H$ and $p$}{n, H and p}}
\label{sec:identity}

On the king lattice the minimum defect over animals of $n$ cells with
bounding-box height $H$ is exactly
\[
  k_{\min}(n,H) \;=\; \left\lceil \frac{n-H}{H-1} \right\rceil
  \qquad (H \ge 2),
\]
with zero violations and equality \emph{attained} in all $672$ cells $(n,H)$
present at $k \le 5$, $n \le 40$ --- so it is sharp.
Equivalently $H \ge (n+k)/(k+1)$, that is $p \le 4n + 4 - \lceil(n-H)/(H-1)\rceil$,
and by transposition the same with the width. Hence $k = 0$ forces $H = W = n$
(the diagonal sticks) and $k = 1$ forces both dimensions at least $(n+1)/2$.

\paragraph{On the square lattice there is essentially nothing.}
$k_{\min}(n,H) = 0$ for $H \in \{1, n\}$ and $1$ for every $H$ between. The
reason is structural: the two square defect-$0$ animals are the horizontal
\emph{and} vertical sticks, at opposite ends of the height range, so no height is
ever more than one defect unit from optimal. Both king defect-$0$ animals are
diagonal sticks with $H = n$, which is why the king bound bites.

\section{The \texorpdfstring{$k = 6$}{k = 6} verdict}
\label{sec:k6}

Four of the statements above were predictions about a value nobody had
computed. The two censuses that decide them (both lattices, $n = 78$,
$k \le 6$) ran on 80-core hardware in 456 shards, merged only on every shard
reporting success, and are \repo{results/perimdefect\_square8\_n78\_k6.txt} and
\repo{results/perimdefect\_square4\_n78\_k6.txt}: $3.95$ and $3.17$ trillion
animals, $42$ and $23$ hours of wall time at $76$-way parallelism. The target
$n = 78$ was chosen because a period-$6$ degree-$6$ fit with two holdouts per
residue class needs $54$ points above onset $24$.

Three of the four hold and one fails.

\begin{enumerate}
\item \textbf{$\mathrm{onset}(6) = 24$ --- confirmed}, and by the scan rather
  than by a fit. A closed form valid from a given $n$ is valid from any larger
  one, so the \emph{smallest} onset at which the division test passes is the
  real one; it is $24$ on both lattices, against the onset formula's $24$.
\item \textbf{The $\Phid_2$ leading diagonal's $15/4$ --- refuted.} The measured
  value is $5/2$. See \S\ref{sec:coefftri}: the closed form is withdrawn, the
  lattice-independence it was riding on survives.
\item \textbf{The $\Phid_3$ exponent $k - 4$ --- confirmed.} $D_6$ carries
  $\Phid_3^2$, and the exponent no longer rests on the single data point at
  $k = 5$.
\item \textbf{$\Phid_4$ absent until $k = 7$ --- confirmed at $k = 6$.} Every
  candidate denominator carrying a $\Phid_4$ also passes the division test, as
  any multiple of the true denominator must; each is then retested with one
  factor removed, and each reduces. The minimal passing denominator is
  $\Phid_1^7\Phid_2^5\Phid_3^2$ on both lattices, with $39$ spare zero
  coefficients.
\end{enumerate}

Item 2 costs less than item 3 failing would have. A wrong \emph{coefficient formula} on the leading diagonal is a
withdrawn table entry. A wrong \emph{denominator} slope would have taken
\S\ref{sec:cyclo}'s $\Phid_d$-first-appears-at-$k = 2d-1$ statement with it, and
with that the reading that each successive periodicity costs two units of
defect. That reading now has two confirmations instead of one.

What $k = 6$ did to this paper: it strengthened the structural claims, which
were predictions about \emph{shape}, and killed the one claim that was an
extrapolated \emph{formula} off three terms. The universality
statement (same period, degree, onset and leading coefficient on both
lattices) now runs through $k = 6$.

\section{Novelty, and what is not ours}
\label{sec:novelty}

A literature-priority pass ran on 2026-08-18
(\repo{docs/priority-passes-2026-08-18.md}). It obtained first the Asinowski
slides and then the full paper behind them, both of which the repository had
recorded as unavailable, and the result is worse than the earlier note here
claimed: the square column reproduces \cite{abz2018}, and three statements
this paper leans on are theirs.

\begin{enumerate}
\item \textbf{The defect identity} $k = 2c + t$ is their $k = e + 2f$
  (Remark~\ref{rem:abzidentity}). It licenses the enumeration prune, so it is
  load-bearing for every measured number in the maximum-end sections.
\item \textbf{Rationality with cyclotomic denominators}, for each fixed defect,
  is their theorem --- stated for polyominoes and for $d$-dimensional polycubes.
  \S\ref{sec:cyclo} is the king analogue of it.
\item Their small cases agree with the square column here: $A(n,2n+2) = 1$ for
  $n \ge 2$, and $A(n,2n+1) = 4(n-2)$ for $n \ge 3$.
\item \textbf{The degree column is their conjecture}, confirmed here and not
  discovered here: their \S4 predicts a highest-degree factor $(x-1)^{k+1}$
  and asymptotics $\gamma n^k$ in any dimension.
\end{enumerate}

A second pass the same day obtained the full text with proofs
(\cite{abz2018}, the polycube companion to the note the slides were drawn
from), which sharpens the verdict in one direction that matters. Their
rationality proof runs by shrinking cuts to a unique reduced representative per
pattern class and bounding the class count through excess cells, $L$-cells and
degree-1 cells; every step of it is stated under face connectivity, down to a
handshake bound at maximum degree $6$. \textbf{The king column is therefore a
separate derivation and not a corollary of their framework} --- the argument's
shape transports, its statements do not. That is the strongest thing this paper
can claim on the structural side, and it is a claim about their proof.

\paragraph{The square column's reach.} \cite{barequetmagal2023}
give an algorithm that emits the closed formulae and generating functions per
defect on the square lattice up to $k = 5$, where $k \le 3$ was previously
published. We hold that paper only by its abstract, so the comparison cannot be
made line by line, but the square column here should be read as agreeing with
it rather than extending it below $k = 6$.

\paragraph{What the pass did not find anywhere.} The king column throughout; the
onset formula $\binom{k+1}{2} + 3$; the two-lattice universality of period,
degree, onset and leading coefficient, now tested through $k = 6$; the
coefficient triangle and its one-diagonal-deep lattice independence; the whole
minimum end, including the attainability criterion, the $P(x)^c$ hull
factorisation and the identification of the diamond tip series with Andrews'
$\phi_2$; and the $n/H/p$ identity. The published minimal-perimeter literature
is structural rather than enumerative (inflation of minimal-perimeter animals
and the constant-isomer conjecture, on the square and hexagonal lattices), and
does not count by perimeter in the way \S\ref{sec:min} does.

\paragraph{What the contribution is.} A published theorem about one
lattice, re-proved for another where the constants differ, plus a family of
structural statements about the pair that the published work does not address
because it does not consider the pair. That is worth writing down and it is not
what an unqualified ``new'' would mean.

\section{Open problems}

\begin{openproblem}[prove the degree and the onset]
Table~\ref{tab:maxend} is interpolation, and its degree column is a published
conjecture (\cite{abz2018}, \S4) confirmed rather than proved. The height
grading has a proof
(polynomial times exponential, degree $\le k$, sharp onset $2k+1$) because a
single-cell row is a cut and the decomposition is exact. The perimeter defect
supports no such cut, which is why the row-transfer approach fails here, so a
proof would need a different mechanism.
\end{openproblem}

\begin{openproblem}[is the two-variable generating function rational?]
$\deg\widetilde N_k = \deg D_k - 1$ exactly, measured on both lattices at
every $k = 2..6$ ($2, 5, 7, 11, 15$) and so growing linearly, as a
rational $F(x,y) = \sum_k G_k(x)y^k$ would require. Seven terms across two
denominator regimes are still not a proof of rationality.
\end{openproblem}

\begin{openproblem}[the hull factorisation]
The $P(x)^c$ argument assumes the corners act independently. That is confirmed
for the 4-corner box, the 6-corner hexagon and the 4-tip diamond. The
$6$-corner/$8$-corner hull factorisation in general is unresolved.
\end{openproblem}

\section{Reproduction}
\label{sec:repro}

\begin{table}[H]
\centering\small
\begin{tabular}{@{}l R{0.56\textwidth}@{}}
\toprule
claim & check \\
\midrule
the enumerator and its gate & \texttt{make build/perimeter\_defect}; \repo{scripts/perimeter\_defect\_gate.sh} \\
Table~\ref{tab:maxend} & \repo{experiments/perimeter\_defect\_fit.py} \\
the onset formula & \repo{experiments/perimeter\_max\_structure.py} \\
the denominators, predicted then divided & \repo{experiments/perimeter\_defect\_gf.py} \texttt{--kmax 6} \\
the onset and $D_k$ minimality, by scan & \repo{experiments/perimeter\_defect\_denominator.py} \texttt{--k 6} \\
$\deg\widetilde N_k = \deg D_k - 1$ & \repo{experiments/perimeter\_defect\_tail\_degree.py} \\
the coefficient triangle & \repo{experiments/perimeter\_defect\_gf.py} \texttt{--compare} \\
the min-end census & \texttt{make build/perimeter\_min}; \repo{scripts/perimeter\_min\_gate.sh} \\
the sharded resumable driver & \repo{scripts/perimeter\_min\_shard\_gate.sh} \\
attainability & \repo{experiments/perimeter\_min\_attainability.py} \\
the tip series & \repo{experiments/cone\_order\_ideals.py} \\
the diamond series, second source & \repo{experiments/diamond\_free\_removals.py} \\
the $n$/$H$/$p$ identity & \repo{results/king\_joint\_nhp\_n9.txt}, \repo{results/king\_triple\_classes\_n9.txt} \\
census data, max end & \repo{results/perimdefect\_square4\_n78\_k6.txt}, \repo{results/perimdefect\_square8\_n78\_k6.txt} and siblings \\
census data, min end & \repo{results/perimmin\_square8\_p48\_r6.txt} and siblings \\
\bottomrule
\end{tabular}
\caption{Where each claim is checked. The brute-force cross-checks that validate
the pruned enumerator are \repo{results/siteperim\_square4\_n20.txt} and
\repo{results/siteperim\_square8\_n14.txt}.}
\label{tab:repro}
\end{table}

\bibliographystyle{plain}
\bibliography{shared/refs}

\end{document}
